\chapter{力与惯性原理}\label{ch:g10:inertia}

停止蹬踏，自行车便滑行直至停下；停止拉动，
雪橇便静止不动。经验告诉我们，运动必须靠
施力来维持 --- 两千年来，物理学也一直这样认为。本章
建立力的概念：力是以牛顿为单位度量的箭头，并阐述力学的第一
条伟大定律：若物体真正不受其他作用，运动物体
将永远保持匀速直线运动。

\section{用力刻画作用}

\begin{definition}[力]\label{def:g10:inertia:force}
\emph{力}\index{力}刻画一种机械作用（推、
拉、吸引）作用于物体上。它由以下特征刻画：
\begin{itemize}
  \item 其\emph{作用点}\index{作用点}；
  \item 方向 --- 作用所沿的直线以及
        沿该直线的指向；
  \item 大小，以\emph{牛顿}\index{牛顿}为单位度量
        （符号 \unit{N}）。
\end{itemize}
力用箭头 $\vect F$ 表示（矢量，同数学
分册），从作用点出发，长度与大小成正比，
按规定比例尺画出。
\end{definition}

\begin{example}[书包的重力]\label{ex:g10:inertia:schoolbag}
一个质量为 $m = \qty{4.0}{kg}$ 的书包，受地球
竖直向下的拉力，即重力（见\cref{ch:g10:universal-gravitation}）：
$P = mg = 4.0 \times 9.81 \approx \qty{39}{N}$ --- 一牛顿大约
相当于一个苹果的重力。取比例 $\qty{1}{cm} \leftrightarrow
\qty{10}{N}$：竖直箭头长 \qty{3.9}{cm}，作用于
书包的重心，方向竖直向下。
\end{example}

\section{力的清单}

\begin{definition}[接触力与非接触力]\label{def:g10:inertia:contact}
\emph{接触力}\index{接触力}作用于两物体相互接触的
地方。\emph{非接触力}\index{非接触力}无需接触：
重力 --- 即
\cref{ch:g10:universal-gravitation}中的万有引力 --- 是本章
唯一使用的非接触力。
\end{definition}

\begin{definition}[常见的接触力]\label{def:g10:inertia:inventory}
四种接触力可涵盖大多数日常情形：
\begin{itemize}
  \item 支持物的\emph{法向力}\index{法向力} $\vect R$，
        垂直于其表面、背离支持物；
  \item 细线、绳索或链条的\emph{张力}\index{张力} $\vect T$，
        沿其方向，将物体拉向绳子一侧；
  \item \emph{摩擦力}\index{摩擦力} $\vect f$，来自
        接触面、空气或水，阻碍滑动或运动；
  \item 发动机、喷气装置或手的\emph{推力}\index{推力}，
        沿施力方向推动物体。
\end{itemize}
\end{definition}

\begin{omfigure}
\begin{tikzpicture}[scale=0.9, font=\small]
  \draw[fill=gray!15] (-2.4,0) rectangle (2.4,-0.25);
  \draw (-2.0,-0.25) -- (-2.0,-1.5);
  \draw (2.0,-0.25) -- (2.0,-1.5);
  \draw[fill=omExo!25] (-0.9,0) rectangle (0.9,0.5);
  \coordinate (G) at (0,0.25);
  \fill (G) circle (1.5pt);
  \draw[-Stealth, omThm, very thick] (G) -- (0,-1.3)
    node[right] {$\vect P$};
  \draw[-Stealth, omDef, very thick] (G) -- (0,1.8)
    node[right] {$\vect R$};
\end{tikzpicture}

{\small 桌上的书：重力 $\vect P$ 与法向力
$\vect R$ --- 同一直线、长度相等、方向相反。}
\end{omfigure}

\begin{example}[被拖雪橇的受力清单]\label{ex:g10:inertia:sled}
雪地上被拖行的雪橇：重力 $\vect P$（非接触力），雪面的
法向力 $\vect R$，绳子的张力 $\vect T$，摩擦力
$\vect f$ 阻碍滑动。除此之外，没有其他物体接触雪橇，
受力清单即完整 --- \emph{重力始终计入，之后每个接触
物体各计一个接触力}。
\end{example}

\section{惯性原理}

\begin{definition}[匀速直线运动]\label{def:g10:inertia:uniform}
若物体在某参考系中做\emph{匀速直线运动}\index{匀速直线运动}，
即其轨迹为一条直线，且以
恒定速率运动：其速度矢量 $\vect v$ 在每一
时刻都相同。静止是 $\vect v = \vect 0$ 的特殊情形。
\end{definition}

\begin{definition}[相互抵消的力]\label{def:g10:inertia:compensate}
作用于物体的力若\emph{相互抵消}\index{相互抵消的力}，
即其合效应为零：首尾相接，各箭头回到
起点。两个力恰好相互抵消，当且仅当它们的作用线
相同、大小相等、方向相反。
\end{definition}

\begin{theorem}[惯性原理]\label{thm:g10:inertia:principle}
\index{惯性原理}\index{惯性}
在地面参考系中，不受力，或所受的力
相互抵消，则物体保持静止或保持匀速直线运动；
反之，静止或做匀速直线运动的物体，所受的力
相互抵消（或不受力）。
\end{theorem}
\admitted

\begin{remark}
没有任何实验能完全隔离一个物体，因此该原理是一条公设
--- 由伽利略窥见，由牛顿明确表述。后文将其升级为一条
联系力与速度\emph{变化}的定律；本科一年级分册
要求判断它在哪些参考系中真正成立。
\end{remark}

\begin{proposition}[由力读出运动，反之亦然]\label{prop:g10:inertia:converse}
在地面参考系中：
\begin{enumerate}
  \item 若物体速度的大小\emph{或}方向发生变化 ---
        作用于它的力\emph{不}相互抵消；
  \item 若物体静止或做匀速直线运动，则其
        所列各力必须相互抵消 --- 据此确定未知量。
\end{enumerate}
\end{proposition}

\begin{proof}
第一条是
\cref{thm:g10:inertia:principle}的逆否命题，第二条则是其逆命题
在完整受力清单上的应用。
\end{proof}

\begin{omfigure}
\begin{tikzpicture}[scale=0.85, font=\small]
  \draw[gray!60] (-0.5,0) -- (10.3,0) node[below left, gray] {冰面};
  \foreach \x in {0.5,2.5,4.5,6.5,8.5} {
    \fill (\x,0.13) ellipse (0.3 and 0.13);
    \draw[-Stealth, omDef, thick] (\x,0.5) -- (\x+1.3,0.5);
  }
  \node[omDef] at (1.15,0.9) {$\vect v$};
  \node[omDef] at (5.4,1.5) {相同时间间隔内步距相等};
\end{tikzpicture}

{\small 冰球脱手后，在相等时间间隔内的位置，
每处各画一个相同的速度矢量 --- 匀速直线运动。}
\end{omfigure}

\begin{example}[冰球]\label{ex:g10:inertia:puck}
冰球被球杆击出后，在光滑冰面上滑过 \qty{30}{m}，
用时 \qty{2.5}{s}，沿直线以恒定速度 $v = 30/2.5 = \qty{12}{m/s}$ 运动。
受力清单：重力与法向力（摩擦力可忽略）。二者
相互抵消，该原理恰好预言了所观察到的运动：
没有任何力推动冰球向前 --- 也无需如此。
\end{example}

\begin{remark}[伽利略与亚里士多德]
亚里士多德认为，每一种运动都需要推动者，日常生活 ---
每次接触都有摩擦力 --- 似乎印证此说。伽利略将摩擦力理想化
即忽略不计：表面越光滑，
越小的推力即可\emph{维持}运动。运动的持续而非静止，才是自然状态。
\end{remark}

\begin{remark}[急刹车时前倾的乘客]
公交车急刹时，站立的乘客向前冲。没有任何力
将他们抛出。由于惯性，乘客的身体保持速度，
而车身的速度减小。安全带提供了受力清单中
所缺的向后的力。
\end{remark}

\begin{omfigure}
\begin{tikzpicture}[scale=0.62, font=\small]
  \draw[gray!60] (-0.3,0) -- (7.3,0);
  \foreach \x in {0.7,2.5,4.3,6.1}
    \fill[omDef] (\x,0.22) ellipse (0.34 and 0.22);
  \draw[-Stealth, omThm, thick] (0.7,0.22) -- (0.7,-1.0)
    node[right] {$\vect P$};
  \draw[-Stealth, omDef, thick] (0.7,0.44) -- (0.7,1.65)
    node[right] {$\vect R$};
  \node at (3.6,-1.7) {力相互抵消：步距相等};
\end{tikzpicture}
\qquad
\begin{tikzpicture}[scale=0.62, font=\small]
  \draw[gray!60] (-0.3,0) -- (7.3,0);
  \foreach \x in {0.7,2.4,3.8,4.9,5.7,6.2}
    \fill[omProp] (\x,0.22) ellipse (0.34 and 0.22);
  \draw[-Stealth, omThm, thick] (0.7,0.22) -- (0.7,-1.0)
    node[right] {$\vect P$};
  \draw[-Stealth, omDef, thick] (0.7,0.44) -- (0.7,1.65)
    node[right] {$\vect R$};
  \draw[-Stealth, omProp, very thick] (0.7,0.22) -- (-0.3,0.22)
    node[above] {$\vect f$};
  \node at (3.6,-1.7) {摩擦力未抵消：冰壶减速};
\end{tikzpicture}

{\small 冰壶在相等时间间隔内的位置。左图：$\vect P$ 与
$\vect R$ 相互抵消 --- 步距相等。右图：未抵消的摩擦力
$\vect f$ 仍存在 --- 步距变小，速度随之改变。}
\end{omfigure}

\section{平衡力：由受力图推断静力学}

\begin{definition}[受力图]\label{def:g10:inertia:diagram}
物体的\emph{受力图}\index{受力图}只画出该物体本身，
将作用\emph{于}该物体的每个力，按比例自其
作用点画出；\emph{由}该物体施加的力从不画入。
\end{definition}

\begin{method}[分析静止情形]\label{met:g10:inertia:method}
\begin{enumerate}
  \item 选定研究对象；在地面参考系中分析。
  \item 列出受力清单：总包括重力，每个接触处各一个接触力。
  \item 画出受力图。
  \item 物体静止，由
        \cref{thm:g10:inertia:principle}可知各力相互抵消：从图中读出未知量的大小
        （一维时取长度相等，二维时取分量）。
\end{enumerate}
\end{method}

\begin{example}[桌上的书]\label{ex:g10:inertia:book}
一本质量为 \qty{0.50}{kg} 的书放在桌上：重力
$P = 0.50 \times 9.81 \approx \qty{4.9}{N}$ 与法向力
$\vect R$。静止时，$\vect R$ 与 $\vect P$ 相互抵消：竖直向上，
$R = \qty{4.9}{N}$。若再加放一本书，$R$ 随之调整 ---
支持面恰好施以所需大小的力。
\end{example}

\begin{omfigure}
\begin{tikzpicture}[scale=0.9, font=\small]
  \draw[fill=gray!15] (-2.6,2.2) rectangle (2.6,2.45);
  \draw (0,0) -- (-2.2,2.2);
  \draw (0,0) -- (2.2,2.2);
  \draw[dashed, gray] (0,0) -- (0,2.2);
  \draw[gray] (0,1.0) arc (90:135:1.0);
  \node[gray] at (-0.42,1.38) {$45^\circ$};
  \draw[fill=omExo!25] (0,-0.35) circle (0.35);
  \draw[-Stealth, omDef, very thick] (0,0) -- (-1.06,1.06)
    node[left] {$\vect T_1$};
  \draw[-Stealth, omDef, very thick] (0,0) -- (1.06,1.06)
    node[right] {$\vect T_2$};
  \draw[-Stealth, omThm, very thick] (0,0) -- (0,-1.5)
    node[right] {$\vect P$};
\end{tikzpicture}

{\small 悬挂的吊灯：两根绳的张力之水平分量由对称性相互抵消，
其竖直分量共同承担重力。}
\end{omfigure}

\begin{example}[两根细线上的吊灯]\label{ex:g10:inertia:lamp}
一盏质量为 \qty{3.0}{kg} 的吊灯（$P \approx \qty{29.4}{N}$）由两根细线悬挂，
每根细线与竖直方向成 $45^\circ$ 角。由对称性可知两拉力
大小相等（$T_1 = T_2 = T$），其水平分量相互抵消；
竖直分量均为 $T \cos 45^\circ$，共同承担重力：
\[
2\,T \cos 45^\circ = P
\qquad\text{故}\qquad
T = \frac{29.4}{2 \times 0.707} \approx \qty{21}{N}
\]
--- 每根细线承受的拉力超过重力的一半：斜拉
会使张力增大。
\end{example}

\section{练习}

\begin{exercise}[$\star$]\label{exo:g10:inertia:1}
列出下列物体所受的力，并判断是接触力还是非接触力：（a）静止在桌上的书；
（b）挂在细线上的吊灯；（c）在光滑冰面上滑行的冰球；
（d）下落的苹果（不计
空气阻力）。
\end{exercise}

\begin{exercise}[$\star$]\label{exo:g10:inertia:2}
求质量为 \qty{250}{g} 的书和 \qty{60}{kg}
的人所受的重力。若取标度 $\qty{1}{cm} \leftrightarrow \qty{1}{N}$，
书的箭头有多长？
为何此标度对这个人毫无用处？
\end{exercise}

\begin{exercise}[$\star$]\label{exo:g10:inertia:3}
在标度为 $\qty{1}{cm} \leftrightarrow \qty{2}{N}$ 的图上，一个力画成
水平向左、长 \qty{3.5}{cm} 的箭头，作用于点 $A$。
求该力的三要素。
\end{exercise}

\begin{exercise}[$\star$]\label{exo:g10:inertia:4}
一本质量为 \qty{1.2}{kg} 的词典放在书架上。画出它的受力图，
并写出每个力的大小与方向。
\end{exercise}

\begin{exercise}[$\star$]\label{exo:g10:inertia:5}
判断下列说法是否正确，并改正错误的：（a）物体总会停下来，除非有力持续推着它；
（b）匀速直线运动不需要未被抵消的力；
（c）静止的物体不受任何力；
（d）若物体的速度发生改变，
则它所受的力不相互抵消。
\end{exercise}

\begin{exercise}[$\star\star$]\label{exo:g10:inertia:6}
冰球释放后，在 \qty{3.0}{s} 内沿直线匀速滑行了 \qty{24}{m}。
求其速率；它所受的力是否相互抵消？
若在混凝土地面上，同一冰球几米之内便停下，
则受力情况有何变化？由\cref{prop:g10:inertia:converse}
可得出什么结论？
\end{exercise}

\begin{exercise}[$\star\star$]\label{exo:g10:inertia:7}
质量为 \qty{70}{kg} 的乘客站在匀速上升的电梯中，
速率为 \qty{1.5}{m/s}。列出乘客所受的力，它们是否相互抵消？
求地板对乘客的支持力大小。若电梯静止，
结果会有变化吗？
\end{exercise}

\begin{exercise}[$\star\star$]\label{exo:g10:inertia:8}
质量为 \qty{1.8}{kg} 的吊灯静止地挂在一根竖直细线上。
求细线的张力。若将其换成两根
\emph{竖直}的细线并平分负载，
则每根细线的张力为多少？
\end{exercise}

\begin{exercise}[$\star\star$]\label{exo:g10:inertia:9}
公交车急刹车时，站立的乘客向前倾；
公交车以恒定速率过弯时，同一乘客则向外倾。
请解释这两种现象，不要虚构“向前”或“向外”的力。
\end{exercise}

\begin{exercise}[$\star\star$]\label{exo:g10:inertia:10}
总质量为 \qty{85}{kg} 的跳伞者以恒定的
\qty{5.0}{m/s} 的速率竖直下降。画出受力图，
并求伞盖所受的空气阻力。
\end{exercise}

\begin{exercise}[$\star\star$]\label{exo:g10:inertia:11}
一名滑水者被船拖着沿直线匀速前进，水平绳索
施加 \qty{300}{N} 的拉力。水对滑水板的水平阻力
多大（须说明理由）？哪两个竖直方向的力相互抵消？
\end{exercise}

\begin{exercise}[$\star\star\star$]\label{exo:g10:inertia:12}
一盏质量为 \qty{2.4}{kg} 的吊灯，由两根与竖直方向夹角为 $40^\circ$ 的
绳索悬挂（$\cos 40^\circ \approx 0.766$）。求每根绳索的张力。
为何绳索越接近水平，张力越大，
而两根恰好水平的绳索为何永远无法拉住吊灯？
\end{exercise}

\begin{exercise}[$\star\star\star$]\label{exo:g10:inertia:13}
一块冰壶被推出并释放后，一边滑行一边逐渐减速，
最后静止。对每个阶段（推、滑、静止）分别画出
受力图，并说明这些力是否相互抵消，同时引用
\cref{thm:g10:inertia:principle} 或 \cref{prop:g10:inertia:converse}。
\end{exercise}

\begin{exercise}[$\star\star\star$]\label{exo:g10:inertia:14}
探测器\emph{旅行者 1 号}在星际空间中无动力滑行，
速率为 \qty{17}{km/s}，发动机已关闭。它为何不需燃料即可保持
其速率？亚里士多德会作出怎样的预测？它一年中行进
多远？先以 \unit{km} 为单位，再以天文单位（$\qty{1}{au} \approx
\qty{1.496e8}{km}$）表示。为何它为改变方向而
\emph{不}改变速率，也须点火推进？
\end{exercise}

\begin{exercise}[$\star\star\star$]\label{exo:g10:inertia:15}
三根绳索拉着一个静止的圆环。在一张网格纸上，每一小格代表
\qty{1}{N}，两个力分别读作 $\vect F_1\,(3, 2)$ 和
$\vect F_2\,(-1, 2)$。求第三个力的分量、大小和
方向。
\end{exercise}

\section{周末问题：三幕中的惯性}

\begin{problem}[{周末问题 --- 机场步道、冰场与降落伞：恒定速度
是受力平衡的标志，以及降落伞为何减小的不是阻力，
而是速率}]
\label{pb:g10:inertia:1}
三个场景 --- 移动步道上的行李箱、冰场上的冰球，
傍晚空中的跳伞者 --- 同一条定律贯穿三者。

\medskip
\noindent\textbf{第 I 部分 --- 机场步道。}
一个 \qty{23}{kg} 的行李箱放在一条以恒定
速率 \qty{0.75}{m/s} 运动的步道上。
\begin{enumerate}
  \item 以地面参考系看，行李箱做怎样的运动？列出它所受的力。
  \item 这些力是否相互抵消？说明理由，并计算两个力的大小。
  \item 将它的受力图与同一行李箱放在地面上时的受力图相比较，
        这一比较说明了怎样的深刻结论？
  \item 步道突然停止，描述行李箱的运动，
        且不得臆造一个向前的力。
  \item 行李箱的主人相对步道以 \qty{1.2}{m/s} 的速率行走。
        她在地面参考系中的速率是多少？\emph{她}所受的力是否相互抵消？
\end{enumerate}

\medskip
\noindent\textbf{第 II 部分 --- 冰场。}
\begin{enumerate}[resume]
  \item 一个 \qty{160}{g} 的冰球以恒定
        \qty{8.0}{m/s} 的速率沿直线滑行。画出其受力图，并标明各力的大小。
  \item 亚里士多德认为冰球需要一个推动者。冰场会如何回应，
        为何日常生活中的表面似乎支持他？
  \item 冰球滑过一段粗糙区域并减速。此处的受力结论是什么？
        哪一个力是罪魁祸首？
  \item 一名滑冰者以恒定速率平滑转弯。她所受的力
        是否相互抵消？注意：速度是哪一类量？
  \item 将受力图与运动对应：(a) 受力相互抵消；(b) 一个额外
        的力与 $\vect v$ 反向；(c) 一个额外的力
        与 $\vect v$ 垂直 --- (i) 直线减速；(ii) 以恒定
        速率转弯；(iii) 匀速直线运动。
\end{enumerate}

\medskip
\noindent\textbf{第 III 部分 --- 降落伞。}
跳伞者连同装备的总质量为 \qty{80}{kg}。
\begin{enumerate}[resume]
  \item 刚离开飞机时，竖直方向的速率与阻力仍为零。
        计算各力；它们是否相互抵消？接下来必须发生什么？
  \item 阻力随速率增大。描述阻力 $<$ 重力时的下落过程。
  \item 降落伞张开之前，跳伞者达到稳定的
        \qty{50}{m/s} 的“终端速度”。求阻力。此时是何种运动？
        这为何是“伪装的惯性”？
  \item 降落伞张开：此时阻力远大于重力，
        速度会如何变化？（跳伞者始终\emph{向下}
        运动。）
  \item 达到新的稳定速率 \qty{5.0}{m/s}。求阻力，
        并与第 13 题比较。
  \item 点睛之问：降落伞改变的是终端阻力 ---
        还是阻力达到重力时的速率？
\end{enumerate}

\medskip
\noindent\textbf{第 IV 部分 --- 结论。}
\begin{enumerate}[resume]
  \item 补全并说明理由：“静止或做直线匀速运动
        $\iff$ \dots”。为何静止不需要单独的定律？
  \item 在哪一幕中，亚里士多德的“运动需要推动者”会
        最明显地失效，伽利略会指向什么？
  \item 一份考卷声称：“跳伞者以恒定
        速率下落，所以他不受力的作用。”请用一句话更正。
  \item 总结：对每一幕，指出相互抵消的那一对力，
        并陈述本题中反复用到二十次的那条定律。
\end{enumerate}
\end{problem}
